\section{Claims of the longer version that do not hold}\label{app:earlier}

The present paper condenses a longer, openly developed text~\cite{companion}.  Many statements there were made before they were audited,
and several do not hold.  They are listed here so that a reader of the longer version can place them.

\paragraph{The Riemann source.}  The longer version began from a separate paper claiming a proof of the Riemann hypothesis and used it as
the clock of the framework.  That proof does not hold, and it is no longer used as a source.  The following statements lost their
justification and are open: the unitary equivalence of the three kernels and the unity of the three objects; the consistency with Bell
inequalities and the reading of string theory and loop quantum gravity as the two compensator sectors, in so far as they used that
equivalence; the Einstein equation from a trace formula, which is not a defined action; the existence and self-adjointness of a Dirac
operator on the antisymmetric sector.  The Riemann hypothesis as a consequence is withdrawn.

\paragraph{Spacing statistics.}  That Riemann zeros follow GUE statistics is empirical; its derivation from the framework is withdrawn.  The
claim that unfolding in the light frame singles out the scale \(\gamma_L\) is contradicted on the first \(2000\) zeros; the optimum lies at \(1\).

\paragraph{The kernel.}  The Minkowski reading of the Hermitian norm (signature \((1,6)\)), a sector rule by the sign of \(\operatorname{Re}z\)
(not Lorentz-invariant), the annotation of \(S^7\) as a group containing \(G_2\), a cascade through all Hopf fibrations as a composition of
Hopf maps, and a conjugate colour triplet on \(\operatorname{span}(e_1,e_2,e_3)\) were wrong and are replaced by the statements of
Sections~\ref{sec:kernel}--\ref{sec:generation}.  A spectral statement about the compensator matrix lacked the hypotheses of
Proposition~\ref{prop:spectral}.

\paragraph{The cosmological constant.}  An earlier argument set \(\Lambda=0\) in the light frame and explained the observed value as a frame
effect.  The trace formula it used is not an action; the value \(\beta_{\rm sym}^2/\ell_P^2\) exceeds the measured \(\Lambda\) by about \(10^{121}\);
one corollary gave two incompatible values for the same quantity; and two remarks placed \(\beta_{\rm sym}^2\) on \(\Lambda\) where a theorem placed
it on \(\Omega_M\).  Only \eqref{eq:omega}, as a present-epoch statement, remains.

\paragraph{Frame-shifted constants.}  A light-frame Planck length \(\gamma_L^{-3/2}\ell_P\), a drift of \(\alpha\) as a frame effect and a vacuum
ratio \(\gamma_L^5\) assumed that the frames differ in \(\hbar\) and \(G_N\).  Under the counting reading they differ in \(\gamma\), and these statements
are not supported.

\paragraph{Six earlier predictions.}  None is confirmed: an improvement of GUE statistics by light-frame unfolding (not confirmed); a Koide
phase \(\theta_K=\pi-\sqrt2\arcsin\beta_{\rm sym}\) (excluded by the lepton masses); a formula \(\beta_{\rm sym}=3^{-1/2}\exp(-\gamma_1/4\pi\gamma_L^5)\)
(a coincidence with an exponent chosen after the fact); \(\Lambda\) as a frame artefact (contradicted); an annual ratio of \(\alpha\)- and
\(G_N\)-drift (corrected, not testable); a Hawking temperature in the tachyon frame (untested).

\paragraph{Readings that went too far.}  A relation presented as a finding held by definition, and the statement that lepton data select the
``corrected'' horizon counting was reversed once the later literature was taken into account (Section~\ref{sec:pheno}).
